\RequirePackage{luatex85}
\documentclass[border=10pt,12pt]{standalone}

\usepackage{cfr-lm}
\usepackage{amssymb}
\usepackage{pgf}
\usepackage{tikz}
\usetikzlibrary{arrows.meta,calc,positioning}

\begin{document}

\sffamily\bfseries

%% Actor x-positions (cm)
%% P=0  S=9  M=19
\def\px{0}
\def\sx{9}
\def\mx{19}
\def\ybot{33.6}

\begin{tikzpicture}[
  >=Stealth,
  yscale=-1,
]

\input{common.tex}

\tikzstyle{actor}=[rectangle, minimum width=3.6cm, minimum height=0.7cm,
  inner sep=4pt, rounded corners=2pt, align=center, font=\bfseries]
\tikzstyle{ll}=[densely dashed, gray!60, line width=0.5pt]
\tikzstyle{msg}=[->, line width=0.65pt]
\tikzstyle{rsp}=[->, densely dashed, line width=0.65pt]
\tikzstyle{broken}=[red!70, densely dotted, line width=0.8pt]
\tikzstyle{lbl}=[font=\small, inner sep=2pt]
\tikzstyle{note}=[rectangle, rounded corners=2pt, inner sep=5pt,
  font=\small, align=center]
\tikzstyle{seclbl}=[font=\small\itshape, text=gray!70!black, anchor=west]
\tikzstyle{selfloop}=[->, line width=0.65pt]

%% ── actor headers ───────────────────────────────────────────────────────
\node[actor, fill=pbox, text=ptxt] (Ph) at (\px,0) {Primary};
\node[actor, fill=sbox, text=stxt] (Sh) at (\sx,0) {Standby};
\node[actor, fill=mbox, text=mtxt] (Mh) at (\mx,0) {Monitor};

%% lifelines
\draw[ll] (\px,0.35) -- (\px,\ybot);
\draw[ll] (\sx,0.35) -- (\sx,\ybot);
\draw[ll] (\mx,0.35) -- (\mx,\ybot);

%% ── steady state ─────────────────────────────────────────────────────────
\node[seclbl] at (-3.4,1.3) {steady state, both on timeline 1};

\draw[msg] (\px,1.6) -- (\sx,1.6)
  node[lbl,midway,above] {WAL stream};
\draw[msg] (\sx,2.6) -- (\mx,2.6);
\node[lbl, above] at (\sx+3.3,2.6) {node\_active(reported=secondary, tli=1)};
\draw[rsp] (\mx,3.2) -- (\sx,3.2);
\node[lbl, above] at (\mx-2.2,3.2) {goal=secondary};

%% ── network partition ────────────────────────────────────────────────────
\node[seclbl] at (-3.4,4.6) {network partition isolates the standby};

\draw[broken] (\sx,5.5) -- (\px,5.5);
\node[red!70!black, font=\Large] at (5.3,5.5) {$\times$};
\draw[broken] (\sx,6.1) -- (\mx,6.1);
\node[red!70!black, font=\Large] at (14.3,6.1) {$\times$};
\node[lbl, above, text=red!70!black] at (\sx,6.8)
  {WAL stream and monitor connection both cut};

%% ── out-of-band promotion ────────────────────────────────────────────────
\node[seclbl] at (-3.4,8.2) {promoted directly at the Postgres level};

\draw[selfloop] (\sx,9.1) .. controls (\sx+2.6,9.1) and (\sx+2.6,10.3) .. (\sx,10.3);
\node[lbl, right] at (\sx+2.7,9.7) {\texttt{pg\_ctl promote}};

\node[note, fill=red!12, draw=red!50, text width=13.5cm, font=\small]
  at (9.5, 11.6)
  {bypasses \texttt{pg\_autoctl} entirely -- the keeper's own \texttt{current\_role}
   assumption (secondary, should be in recovery) is about to stop matching
   Postgres reality (promoted, writable)};

\draw[selfloop] (\sx,13.0) .. controls (\sx+2.6,13.0) and (\sx+2.6,14.0) .. (\sx,14.0);
\node[lbl, right] at (\sx+2.7,13.5) {local writes + \texttt{CHECKPOINT}};

\node[note, fill=gray!12, draw=gray!55, text width=11cm, font=\small]
  at (9.5, 15.2)
  {standby is now on a genuine local timeline 2 -- WAL no other node has,
   and no other node can ever stream};

%% ── reconnects, looks routine ────────────────────────────────────────────
\node[seclbl] at (-3.4,16.6) {network restored -- reconnects, still ``secondary''};

\draw[msg] (\sx,17.5) -- (\mx,17.5);
\node[lbl, above] at (\sx+3.3,17.5) {node\_active(reported=secondary, tli=2)};
\draw[rsp] (\mx,18.1) -- (\sx,18.1);
\node[lbl, above] at (\mx-2.2,18.1) {goal=secondary};

\node[note, fill=async, draw=gray!50, text width=15cm, font=\small]
  at (9.5, 19.6)
  {no sibling standby to disagree with, so auto-detection reads this as a
   legitimate lineage rather than a fork. The divergence is durably
   published to \texttt{pgautofailover.node\_timeline\_history} but stays
   invisible until an operator pins the correct timeline};

%% ── operator pins ground truth ───────────────────────────────────────────
\node[seclbl] at (-3.4,21.4) {operator pins ground truth};

\draw[selfloop] (\mx,21.6) .. controls (\mx+2.6,21.6) and (\mx+2.6,22.6) .. (\mx,22.6);
\node[lbl, right] at (\mx+2.7,22.1) {\texttt{pg\_autoctl accept timeline --tli 1}};

\draw[msg] (\mx,23.5) -- (\sx,23.5);
\node[lbl, above] at (\mx-2.2,23.5) {goal=catchingup};

\node[note, fill=gray!12, draw=gray!55, text width=13cm, font=\small]
  at (9.5, 24.7)
  {the monitor re-checks every currently-secondary node's ancestry against
   the freshly-pinned lineage right away -- pushed to catchingup within
   about a second, no maintenance toggle and no incidental health-check
   cycle needed};

%% ── ancestry check ───────────────────────────────────────────────────────
\node[seclbl] at (-3.4,26.1) {ancestry check};

\draw[selfloop] (\sx,27.0) .. controls (\sx+2.6,27.0) and (\sx+2.6,28.0) .. (\sx,28.0);
\node[lbl, right] at (\sx+2.7,27.5) {walk primary's real \texttt{TIMELINE\_HISTORY}};

\node[note, fill=red!12, draw=red!50, text width=13cm, font=\small]
  at (9.5, 29.2)
  {tli 2 is not an ancestor of tli 1 -- genuine divergence, not lag.
   Pre-fix: Postgres refuses the reconnect
   (\emph{requested timeline 2 is not a child of this server's history})
   and the standby retries the same doomed reconnect forever};

%% ── pg_rewind, rejoins ───────────────────────────────────────────────────
\node[seclbl] at (-3.4,30.6) {\texttt{pg\_rewind}, rejoins cleanly};

\draw[selfloop] (\sx,31.5) .. controls (\sx+2.6,31.5) and (\sx+2.6,32.3) .. (\sx,32.3);
\node[lbl, right] at (\sx+2.7,31.9) {\texttt{pg\_rewind} onto tli 1 (or fresh
  \texttt{pg\_basebackup} if unreachable)};

\node[note, fill=green!15, draw=green!50!black,
      font=\small\bfseries, minimum width=9cm]
  at (4.5, 33.3) {$\checkmark$~~rejoins as secondary on tli 1, divergent WAL discarded};

%% ── actor footers ───────────────────────────────────────────────────────
\node[actor, fill=pbox, text=ptxt] at (\px, \ybot+0.35) {Primary};
\node[actor, fill=sbox, text=stxt] at (\sx, \ybot+0.35) {Standby};
\node[actor, fill=mbox, text=mtxt] at (\mx, \ybot+0.35) {Monitor};

\end{tikzpicture}
\end{document}
